\documentclass[11pt]{article}
\usepackage{estilo_practica}
\begin{document}
\hypersetup{pageanchor=false}
\begin{titlepage}
\thispagestyle{empty}
{\sffamily\color{Tinta}\large Universidad Nacional Autónoma de México\par
\normalsize Facultad de Ciencias\par}
\vspace{9mm}
{\color{Acento}\rule{\linewidth}{1.2pt}\par}
\vspace{14mm}
{\sffamily\color{Acento}\large FUNDAMENTOS DE BASES DE DATOS\par}
\vspace{6mm}
{\sffamily\color{Tinta}\fontsize{32}{38}\selectfont\bfseries Práctica 05\par}
\vspace{4mm}
{\sffamily\color{Tinta}\fontsize{21}{27}\selectfont Lenguaje de Definición de Datos\par}
\vspace{7mm}
{\large Identificación y contenido de la entrega\par}
\vfill
{\sffamily\color{Tinta}\large\bfseries Equipo ÑIPGG\par}
\vspace{4mm}
\begin{tabularx}{\linewidth}{@{}Xr@{}}
\toprule
\textbf{Nombre completo} & \textbf{Número de cuenta}\\\midrule
Omar Alejandro Juárez Larrondo & 322245244\\[5pt]
Yahir León Bautista & 322176542\\[5pt]
Diego Hernández Gómez & 321069942\\[5pt]
Miguel Ángel Jiménez Ramírez & 119000887\\[5pt]
Ana Lilia Carballido Camacateco & 315314601\\\bottomrule
\end{tabularx}\par
\vspace{9mm}
Profesor: Gerardo Avilés Rosas\\
Semestre 2027-I · 26 de octubre de 2026\par
\vspace{7mm}
{\color{Linea}\rule{\linewidth}{.6pt}\par}
\end{titlepage}
\hypersetup{pageanchor=true}

\section*{Contenido de la entrega}
\begin{description}
\item[README\_ÑIPGG.pdf] Identificación del equipo y descripción del contenido.
\item[Docs/Práctica05.pdf] Dominios y restricciones de las 146 columnas de las 30 tablas,
  evidencias de restauración del respaldo y modelos gráficos de referencia.
\item[SQL/DDL.sql] SQL autónomo ensamblado por fases: tablas, dominios,
  claves y referencias. Validado en una base vacía con PostgreSQL 18 y pruebas de integración.
\item[Diagramas/ERÑIPGG.png] Imagen del modelo entidad--relación.
\item[Diagramas/ERÑIPGG.drawio] Fuente editable del modelo entidad--relación.
\item[Diagramas/RelacionalÑIPGG.png] Imagen del modelo relacional.
\item[Diagramas/RelacionalÑIPGG.drawio] Fuente editable del modelo relacional.
\end{description}
El archivo comprimido contiene exclusivamente
los siete archivos enumerados. Las fuentes de trabajo, pruebas y registros
se conservan fuera del paquete.
\end{document}
